%
% Tulisan "Metode" dapat diganti menjadi "Kerangka Pikir"
%
\section{Metode}
%======================================================%
%       TULIS METODE ATAU KERANGKA PIKIR DI SINI       %
%                                                      %
%          Mulailah menulis dari \subsection           %
%      \subsection > \subsubsection > \paragraph       %
%      > \subparagraph                                 %
%                                                      %
% Jangan Lupa untuk Menghapus Contoh Soal di Bawah Ini %
%======================================================%

Bagian metode menjelaskan secara terperinci mengenai pendekatan, metode, teknik, serta prosedur yang digunakan sepanjang proses penelitian. Penulis wajib menguraikan instrumen yang digunakan dalam pengumpulan, pengolahan, hingga penafsiran dan analisis data. Setiap pemilihan metode atau instrumen tersebut harus disertai dengan penjelasan alasan yang logis dan ilmiah mengapa pendekatan tersebut paling cocok untuk memecahkan masalah penelitian. Seluruh tahapan dituliskan mengalir dalam paragraf yang padat untuk menunjukkan bahwa penelitian dilakukan secara sistematis, terukur, dan dapat dipertanggungjawabkan validitasnya.

Bagian kerangka pikir memuat penjelasan serta gambaran komprehensif mengenai teori, konsep, pendapat ahli, dan informasi ilmiah lain yang terkait langsung dengan topik bahasan. Seluruh pustaka yang diulas di sini diaplikasikan sebagai landasan utama serta sumber rujukan untuk membedah, menganalisis, dan menafsirkan permasalahan beserta solusi yang ditawarkan. Narasi disusun secara kritis untuk membangun argumentasi yang kokoh, sehingga pembaca dapat melihat dengan jelas peta konsep dan dasar logika yang digunakan penulis dalam memetakan fenomena yang sedang dikaji.